\section{万历朝的补充编年节点（一）}
这一组条目把此前尚未设导航入口的年度材料接回本章脉络。它们不是把同年材料机械等同为因果，而是在可核查的时间位置上，分别保留行动、行政记录与社会经验之间的距离。

\subsection{1578年前后的材料线索}

\eventanchor{MM-1578-EMPRESS-WANG-INSTALLATION}{明·1578（万历六年）六月}
\paragraph{1578（万历六年）六月：万历帝册立王氏为皇后，完成幼帝成年后的皇后册立仪式} 这一节点应放回本章已经展开的权力、财政、地方社会与边地关系中理解。账本把它出现的条件概括为：万历帝已亲政，皇室需要依礼确定皇后以组织后宫名分、祭祀和未来继嗣的程序。 随后可追踪的变化是：王氏取得皇后名号，后宫礼制与外戚关系有了正式支点；继承问题仍取决于子嗣、册立和宫廷政治。 可由《神宗实录》万历六年六月乙卯条；《明史·孝端显皇后传》；明代后妃册立仪文 互相核对；不过，册立日期、礼仪和诏令可靠；礼制记录突出朝廷秩序，不能据此推知帝后关系、后宫实际权力或后继生育结果。 因而这里标示的是材料能够支持的行动与制度位置，而不是对动机、地方执行或普通人经验的无保留复原。

\eventanchor{MM-1578-NATIONWIDE-SURVEY}{明·1578（万历六年）}
\paragraph{1578（万历六年）：朝廷令各省清丈田亩、核查隐漏，为赋役重整收集册籍} 这一节点应放回本章已经展开的权力、财政、地方社会与边地关系中理解。账本把它出现的条件概括为：旧有登记与实际占有、赋额脱节，中央财政和改革政策需要更可比的田亩、户役信息。 随后可追踪的变化是：各地重编册籍并暴露隐漏、产权和役额争议；清丈为一条鞭的区域推行提供条件，却不决定其统一形式。 可由《神宗实录》万历六年户部条；《明史·食货志》；徽州、江南等地清丈册与契约 互相核对；不过，清丈命令可靠；各省启动、完结与报亩年份不同，中央汇总数不能当作同一年度的全国实测。 因而这里标示的是材料能够支持的行动与制度位置，而不是对动机、地方执行或普通人经验的无保留复原。

\eventanchor{ML-1579-001}{明·1579（万历七年）}
\paragraph{1579（万历七年）：东林运动：私立东林书院全部关闭} 这一节点应放回本章已经展开的权力、财政、地方社会与边地关系中理解。账本把它出现的条件概括为：继承、官僚协调或皇权运作的压力使问题进入朝廷程序。 随后可追踪的变化是：它影响后续人事与政策议程；动机和派系标签不能由单一叙事裁定。 可由《神宗实录》万历七年条；《明史·本纪》及相关列传；诏令、奏疏或会典版本 互相核对；不过，译写候选以编年材料定位；实录记载反映朝廷选择，崇祯期须另以奏疏、地方及敌对政权材料交叉。 因而这里标示的是材料能够支持的行动与制度位置，而不是对动机、地方执行或普通人经验的无保留复原。

\eventanchor{ML-1579-002}{明·1579（万历七年）}
\paragraph{1579（万历七年）：朝鲜战争、海上贸易和朝贡使节的本年多方记录提供跨政权比较，外交语言不能替代地方经验。} 这一节点应放回本章已经展开的权力、财政、地方社会与边地关系中理解。账本把它出现的条件概括为：朝贡名义、边境安全与港口贸易的利益使交涉持续发生。 随后可追踪的变化是：它为后续册封、互市或海防安排留下文书与跨政权比较线索。 可由《神宗实录》万历七年条；《明史·外国传》；《朝鲜王朝实录》、琉球《历代宝案》或海外档案 互相核对；不过，桥接条目只标示可回查的年度问题与材料链；实录有中央选择性，崇祯期尤其需要奏疏、地方、朝鲜及满文材料交叉。 因而这里标示的是材料能够支持的行动与制度位置，而不是对动机、地方执行或普通人经验的无保留复原。

\eventanchor{ML-1579-003}{明·1579（万历七年）}
\paragraph{1579（万历七年）：本年灾荒、河工和地方赈济材料可用于研究风险分配，灾害叙事和统计数字应区分。} 这一节点应放回本章已经展开的权力、财政、地方社会与边地关系中理解。账本把它出现的条件概括为：自然风险与地方报告促成朝廷或地方的应对记录。 随后可追踪的变化是：赈济、迁徙和损失的实际范围仍须以受灾地区材料分别核验。 可由《神宗实录》万历七年条；《明史·五行志》；受灾府县地方志与奏疏 互相核对；不过，桥接条目只标示可回查的年度问题与材料链；实录有中央选择性，崇祯期尤其需要奏疏、地方、朝鲜及满文材料交叉。 因而这里标示的是材料能够支持的行动与制度位置，而不是对动机、地方执行或普通人经验的无保留复原。

\eventanchor{ML-1579-004}{明·1579（万历七年）}
\paragraph{1579（万历七年）：辽东、西北、朝鲜或西南军务的本年调兵与战报记录跨区域动员，补给与兵额需要独立核验。} 这一节点应放回本章已经展开的权力、财政、地方社会与边地关系中理解。账本把它出现的条件概括为：前期边防、地方武装或跨区域资源调动的压力推动了该行动。 随后可追踪的变化是：它改变了后续调兵、补给或地方秩序的条件；战果和伤亡须分开核对。 可由《神宗实录》万历七年条；《明史·兵志》及相关列传；边镇、地方志或对方材料 互相核对；不过，桥接条目只标示可回查的年度问题与材料链；实录有中央选择性，崇祯期尤其需要奏疏、地方、朝鲜及满文材料交叉。 因而这里标示的是材料能够支持的行动与制度位置，而不是对动机、地方执行或普通人经验的无保留复原。

\eventanchor{ML-1579-005}{明·1579（万历七年）}
\paragraph{1579（万历七年）：清丈、一条鞭、漕运或军饷的本年文书构成万历财政的可核查节点，定额、报数和实收必须分列。} 这一节点应放回本章已经展开的权力、财政、地方社会与边地关系中理解。账本把它出现的条件概括为：财政征解、交通工程或货币支付的压力使相关措施进入政务。 随后可追踪的变化是：其法定安排不等于地方执行，后续须以地域材料追踪负担与成效。 可由《神宗实录》万历七年条；《明会典》相关条目；地方志、黄册/契约/河工材料 互相核对；不过，桥接条目只标示可回查的年度问题与材料链；实录有中央选择性，崇祯期尤其需要奏疏、地方、朝鲜及满文材料交叉。 因而这里标示的是材料能够支持的行动与制度位置，而不是对动机、地方执行或普通人经验的无保留复原。

\eventanchor{MM-1579-ACADEMY-RESTRICTIONS}{明·1579（万历七年）}
\paragraph{1579（万历七年）：张居正政府以禁讲学、限制书院活动的诏令约束地方士人结社} 这一节点应放回本章已经展开的权力、财政、地方社会与边地关系中理解。账本把它出现的条件概括为：张居正担忧讲学社群与地方结交影响官员考成和政治批评，中央希望将士人活动重新纳入官学与官僚秩序。 随后可追踪的变化是：部分书院公开活动受限，士人转向家塾、文会或其他交往形式；禁令也为后来的学术自治记忆提供素材。 可由《神宗实录》万历七年书院条；张居正公牍与奏疏；书院碑记及地方志 互相核对；不过，禁令和中央政策方向可靠；各地书院是否停讲、改名或私下延续须以地方志、碑刻和文集逐处核验，不能以诏令写成思想生活完全中断。 因而这里标示的是材料能够支持的行动与制度位置，而不是对动机、地方执行或普通人经验的无保留复原。

\eventanchor{ML-1580-001}{明·1580（万历八年）}
\paragraph{1580（万历八年）：单鞭法：简化税法} 这一节点应放回本章已经展开的权力、财政、地方社会与边地关系中理解。账本把它出现的条件概括为：财政征解、交通工程或货币支付的压力使相关措施进入政务。 随后可追踪的变化是：其法定安排不等于地方执行，后续须以地域材料追踪负担与成效。 可由《神宗实录》万历八年条；《明会典》相关条目；地方志、黄册/契约/河工材料 互相核对；不过，译写候选以编年材料定位；实录记载反映朝廷选择，崇祯期须另以奏疏、地方及敌对政权材料交叉。 因而这里标示的是材料能够支持的行动与制度位置，而不是对动机、地方执行或普通人经验的无保留复原。

\eventanchor{ML-1580-002}{明·1580（万历八年）}
\paragraph{1580（万历八年）：官员们批评万历皇帝的疏忽和个人生活的不当行为} 这一节点应放回本章已经展开的权力、财政、地方社会与边地关系中理解。账本把它出现的条件概括为：继承、官僚协调或皇权运作的压力使问题进入朝廷程序。 随后可追踪的变化是：它影响后续人事与政策议程；动机和派系标签不能由单一叙事裁定。 可由《神宗实录》万历八年条；《明史·本纪》及相关列传；诏令、奏疏或会典版本 互相核对；不过，译写候选以编年材料定位；实录记载反映朝廷选择，崇祯期须另以奏疏、地方及敌对政权材料交叉。 因而这里标示的是材料能够支持的行动与制度位置，而不是对动机、地方执行或普通人经验的无保留复原。

\eventanchor{ML-1580-003}{明·1580（万历八年）}
\paragraph{1580（万历八年）：清丈、一条鞭、漕运或军饷的本年文书构成万历财政的可核查节点，定额、报数和实收必须分列。} 这一节点应放回本章已经展开的权力、财政、地方社会与边地关系中理解。账本把它出现的条件概括为：财政征解、交通工程或货币支付的压力使相关措施进入政务。 随后可追踪的变化是：其法定安排不等于地方执行，后续须以地域材料追踪负担与成效。 可由《神宗实录》万历八年条；《明会典》相关条目；地方志、黄册/契约/河工材料 互相核对；不过，桥接条目只标示可回查的年度问题与材料链；实录有中央选择性，崇祯期尤其需要奏疏、地方、朝鲜及满文材料交叉。 因而这里标示的是材料能够支持的行动与制度位置，而不是对动机、地方执行或普通人经验的无保留复原。

\eventanchor{ML-1580-004}{明·1580（万历八年）}
\paragraph{1580（万历八年）：朝鲜战争、海上贸易和朝贡使节的本年多方记录提供跨政权比较，外交语言不能替代地方经验。} 这一节点应放回本章已经展开的权力、财政、地方社会与边地关系中理解。账本把它出现的条件概括为：朝贡名义、边境安全与港口贸易的利益使交涉持续发生。 随后可追踪的变化是：它为后续册封、互市或海防安排留下文书与跨政权比较线索。 可由《神宗实录》万历八年条；《明史·外国传》；《朝鲜王朝实录》、琉球《历代宝案》或海外档案 互相核对；不过，桥接条目只标示可回查的年度问题与材料链；实录有中央选择性，崇祯期尤其需要奏疏、地方、朝鲜及满文材料交叉。 因而这里标示的是材料能够支持的行动与制度位置，而不是对动机、地方执行或普通人经验的无保留复原。

\eventanchor{ML-1580-005}{明·1580（万历八年）}
\paragraph{1580（万历八年）：考成、人事、立储和留中等本年材料标记皇帝、内阁与言官的协调关系，不可只用党争标签解释。} 这一节点应放回本章已经展开的权力、财政、地方社会与边地关系中理解。账本把它出现的条件概括为：继承、官僚协调或皇权运作的压力使问题进入朝廷程序。 随后可追踪的变化是：它影响后续人事与政策议程；动机和派系标签不能由单一叙事裁定。 可由《神宗实录》万历八年条；《明史·本纪》及相关列传；诏令、奏疏或会典版本 互相核对；不过，桥接条目只标示可回查的年度问题与材料链；实录有中央选择性，崇祯期尤其需要奏疏、地方、朝鲜及满文材料交叉。 因而这里标示的是材料能够支持的行动与制度位置，而不是对动机、地方执行或普通人经验的无保留复原。

\eventanchor{MM-1580-ONE-WHIP-ORDERS}{明·1580（万历八年）}
\paragraph{1580（万历八年）：朝廷以清丈和赋役整顿为基础，推动各地合并田赋、丁役与杂派的征收办法} 这一节点应放回本章已经展开的权力、财政、地方社会与边地关系中理解。账本把它出现的条件概括为：清丈揭示登记与赋役的脱节，役法名目繁复且征解成本高，张居正政府希望简化责任与提高可核算性。 随后可追踪的变化是：更多地区以田亩为主要计征基础并折纳部分银两；不同地方仍保留实物、差役或本地变通。 可由《神宗实录》万历八年户部条；《明史·食货志》；各地赋役全书、清丈册和契约 互相核对；不过，中央推动赋役合并的方向可靠；“一条鞭法”在各省府县的项目、折银比例和实施节奏并不一致。 因而这里标示的是材料能够支持的行动与制度位置，而不是对动机、地方执行或普通人经验的无保留复原。

\eventanchor{ML-1593-003}{明·1593（万历二十一年）}
\paragraph{1593（万历二十一年）：嘉运河扩建开始} 这一节点应放回本章已经展开的权力、财政、地方社会与边地关系中理解。账本把它出现的条件概括为：财政征解、交通工程或货币支付的压力使相关措施进入政务。 随后可追踪的变化是：其法定安排不等于地方执行，后续须以地域材料追踪负担与成效。 可由《神宗实录》万历二十一年条；《明会典》相关条目；地方志、黄册/契约/河工材料 互相核对；不过，译写候选以编年材料定位；实录记载反映朝廷选择，崇祯期须另以奏疏、地方及敌对政权材料交叉。 因而这里标示的是材料能够支持的行动与制度位置，而不是对动机、地方执行或普通人经验的无保留复原。

\eventanchor{ML-1593-004}{明·1593（万历二十一年）}
\paragraph{1593（万历二十一年）：中、初年级作业采用抽签方式分配} 这一节点应放回本章已经展开的权力、财政、地方社会与边地关系中理解。账本把它出现的条件概括为：继承、官僚协调或皇权运作的压力使问题进入朝廷程序。 随后可追踪的变化是：它影响后续人事与政策议程；动机和派系标签不能由单一叙事裁定。 可由《神宗实录》万历二十一年条；《明史·本纪》及相关列传；诏令、奏疏或会典版本 互相核对；不过，译写候选以编年材料定位；实录记载反映朝廷选择，崇祯期须另以奏疏、地方及敌对政权材料交叉。 因而这里标示的是材料能够支持的行动与制度位置，而不是对动机、地方执行或普通人经验的无保留复原。

\eventanchor{ML-1593-005}{明·1593（万历二十一年）}
\paragraph{1593（万历二十一年）：明朝官员每年向中国帆船颁发十张许可证，允许其在台湾北部进行贸易} 这一节点应放回本章已经展开的权力、财政、地方社会与边地关系中理解。账本把它出现的条件概括为：财政征解、交通工程或货币支付的压力使相关措施进入政务。 随后可追踪的变化是：其法定安排不等于地方执行，后续须以地域材料追踪负担与成效。 可由《神宗实录》万历二十一年条；《明会典》相关条目；地方志、黄册/契约/河工材料 互相核对；不过，译写候选以编年材料定位；实录记载反映朝廷选择，崇祯期须另以奏疏、地方及敌对政权材料交叉。 因而这里标示的是材料能够支持的行动与制度位置，而不是对动机、地方执行或普通人经验的无保留复原。

\eventanchor{ML-1594-001}{明·1594（万历二十二年）}
\paragraph{1594（万历二十二年）：亳州之乱：明军在四川大败} 这一节点应放回本章已经展开的权力、财政、地方社会与边地关系中理解。账本把它出现的条件概括为：前期边防、地方武装或跨区域资源调动的压力推动了该行动。 随后可追踪的变化是：它改变了后续调兵、补给或地方秩序的条件；战果和伤亡须分开核对。 可由《神宗实录》万历二十二年条；《明史·兵志》及相关列传；边镇、地方志或对方材料 互相核对；不过，译写候选以编年材料定位；实录记载反映朝廷选择，崇祯期须另以奏疏、地方及敌对政权材料交叉。 因而这里标示的是材料能够支持的行动与制度位置，而不是对动机、地方执行或普通人经验的无保留复原。

\eventanchor{ML-1594-002}{明·1594（万历二十二年）}
\paragraph{1594（万历二十二年）：考成、人事、立储和留中等本年材料标记皇帝、内阁与言官的协调关系，不可只用党争标签解释。} 这一节点应放回本章已经展开的权力、财政、地方社会与边地关系中理解。账本把它出现的条件概括为：继承、官僚协调或皇权运作的压力使问题进入朝廷程序。 随后可追踪的变化是：它影响后续人事与政策议程；动机和派系标签不能由单一叙事裁定。 可由《神宗实录》万历二十二年条；《明史·本纪》及相关列传；诏令、奏疏或会典版本 互相核对；不过，桥接条目只标示可回查的年度问题与材料链；实录有中央选择性，崇祯期尤其需要奏疏、地方、朝鲜及满文材料交叉。 因而这里标示的是材料能够支持的行动与制度位置，而不是对动机、地方执行或普通人经验的无保留复原。

\eventanchor{ML-1594-003}{明·1594（万历二十二年）}
\paragraph{1594（万历二十二年）：辽东、西北、朝鲜或西南军务的本年调兵与战报记录跨区域动员，补给与兵额需要独立核验。} 这一节点应放回本章已经展开的权力、财政、地方社会与边地关系中理解。账本把它出现的条件概括为：前期边防、地方武装或跨区域资源调动的压力推动了该行动。 随后可追踪的变化是：它改变了后续调兵、补给或地方秩序的条件；战果和伤亡须分开核对。 可由《神宗实录》万历二十二年条；《明史·兵志》及相关列传；边镇、地方志或对方材料 互相核对；不过，桥接条目只标示可回查的年度问题与材料链；实录有中央选择性，崇祯期尤其需要奏疏、地方、朝鲜及满文材料交叉。 因而这里标示的是材料能够支持的行动与制度位置，而不是对动机、地方执行或普通人经验的无保留复原。

\subsection{1594年前后的材料线索}

\eventanchor{ML-1594-004}{明·1594（万历二十二年）}
\paragraph{1594（万历二十二年）：本年灾荒、河工和地方赈济材料可用于研究风险分配，灾害叙事和统计数字应区分。} 这一节点应放回本章已经展开的权力、财政、地方社会与边地关系中理解。账本把它出现的条件概括为：自然风险与地方报告促成朝廷或地方的应对记录。 随后可追踪的变化是：赈济、迁徙和损失的实际范围仍须以受灾地区材料分别核验。 可由《神宗实录》万历二十二年条；《明史·五行志》；受灾府县地方志与奏疏 互相核对；不过，桥接条目只标示可回查的年度问题与材料链；实录有中央选择性，崇祯期尤其需要奏疏、地方、朝鲜及满文材料交叉。 因而这里标示的是材料能够支持的行动与制度位置，而不是对动机、地方执行或普通人经验的无保留复原。

\eventanchor{ML-1594-005}{明·1594（万历二十二年）}
\paragraph{1594（万历二十二年）：朝鲜战争、海上贸易和朝贡使节的本年多方记录提供跨政权比较，外交语言不能替代地方经验。} 这一节点应放回本章已经展开的权力、财政、地方社会与边地关系中理解。账本把它出现的条件概括为：朝贡名义、边境安全与港口贸易的利益使交涉持续发生。 随后可追踪的变化是：它为后续册封、互市或海防安排留下文书与跨政权比较线索。 可由《神宗实录》万历二十二年条；《明史·外国传》；《朝鲜王朝实录》、琉球《历代宝案》或海外档案 互相核对；不过，桥接条目只标示可回查的年度问题与材料链；实录有中央选择性，崇祯期尤其需要奏疏、地方、朝鲜及满文材料交叉。 因而这里标示的是材料能够支持的行动与制度位置，而不是对动机、地方执行或普通人经验的无保留复原。

\eventanchor{ML-1595-001}{明·1595（万历二十三年）}
\paragraph{1595（万历二十三年）：清丈、一条鞭、漕运或军饷的本年文书构成万历财政的可核查节点，定额、报数和实收必须分列。（材料核查重点：诏令或奏议的形成与发受链）} 这一节点应放回本章已经展开的权力、财政、地方社会与边地关系中理解。账本把它出现的条件概括为：财政征解、交通工程或货币支付的压力使相关措施进入政务。 随后可追踪的变化是：其法定安排不等于地方执行，后续须以地域材料追踪负担与成效。 可由《神宗实录》万历二十三年条；《明会典》相关条目；地方志、黄册/契约/河工材料 互相核对；不过，桥接条目只标示可回查的年度问题与材料链；实录有中央选择性，崇祯期尤其需要奏疏、地方、朝鲜及满文材料交叉。；本条特别核对诏令或奏议的形成与发受链。 因而这里标示的是材料能够支持的行动与制度位置，而不是对动机、地方执行或普通人经验的无保留复原。

\eventanchor{ML-1595-002}{明·1595（万历二十三年）}
\paragraph{1595（万历二十三年）：辽东、西北、朝鲜或西南军务的本年调兵与战报记录跨区域动员，补给与兵额需要独立核验。} 这一节点应放回本章已经展开的权力、财政、地方社会与边地关系中理解。账本把它出现的条件概括为：前期边防、地方武装或跨区域资源调动的压力推动了该行动。 随后可追踪的变化是：它改变了后续调兵、补给或地方秩序的条件；战果和伤亡须分开核对。 可由《神宗实录》万历二十三年条；《明史·兵志》及相关列传；边镇、地方志或对方材料 互相核对；不过，桥接条目只标示可回查的年度问题与材料链；实录有中央选择性，崇祯期尤其需要奏疏、地方、朝鲜及满文材料交叉。 因而这里标示的是材料能够支持的行动与制度位置，而不是对动机、地方执行或普通人经验的无保留复原。

\eventanchor{ML-1595-003}{明·1595（万历二十三年）}
\paragraph{1595（万历二十三年）：朝鲜战争、海上贸易和朝贡使节的本年多方记录提供跨政权比较，外交语言不能替代地方经验。} 这一节点应放回本章已经展开的权力、财政、地方社会与边地关系中理解。账本把它出现的条件概括为：朝贡名义、边境安全与港口贸易的利益使交涉持续发生。 随后可追踪的变化是：它为后续册封、互市或海防安排留下文书与跨政权比较线索。 可由《神宗实录》万历二十三年条；《明史·外国传》；《朝鲜王朝实录》、琉球《历代宝案》或海外档案 互相核对；不过，桥接条目只标示可回查的年度问题与材料链；实录有中央选择性，崇祯期尤其需要奏疏、地方、朝鲜及满文材料交叉。 因而这里标示的是材料能够支持的行动与制度位置，而不是对动机、地方执行或普通人经验的无保留复原。

\eventanchor{ML-1595-004}{明·1595（万历二十三年）}
\paragraph{1595（万历二十三年）：清丈、一条鞭、漕运或军饷的本年文书构成万历财政的可核查节点，定额、报数和实收必须分列。（材料核查重点：报称信息的来源、抄传与行政分类）} 这一节点应放回本章已经展开的权力、财政、地方社会与边地关系中理解。账本把它出现的条件概括为：财政征解、交通工程或货币支付的压力使相关措施进入政务。 随后可追踪的变化是：其法定安排不等于地方执行，后续须以地域材料追踪负担与成效。 可由《神宗实录》万历二十三年条；《明会典》相关条目；地方志、黄册/契约/河工材料 互相核对；不过，桥接条目只标示可回查的年度问题与材料链；实录有中央选择性，崇祯期尤其需要奏疏、地方、朝鲜及满文材料交叉。；本条特别核对报称信息的来源、抄传与行政分类。 因而这里标示的是材料能够支持的行动与制度位置，而不是对动机、地方执行或普通人经验的无保留复原。

\eventanchor{ML-1595-005}{明·1595（万历二十三年）}
\paragraph{1595（万历二十三年）：书院、科举、出版与宗教活动的本年线索可连接知识网络，存世文本有阶层与地域偏差。} 这一节点应放回本章已经展开的权力、财政、地方社会与边地关系中理解。账本把它出现的条件概括为：制度、教育或知识传播的需要推动了相关行动或文本形成。 随后可追踪的变化是：它提供制度和传播史的锚点，不能据此推断社会整体接受。 可由《神宗实录》万历二十三年条；《明史·选举志》或《艺文志》；文集、碑刻或书目材料 互相核对；不过，桥接条目只标示可回查的年度问题与材料链；实录有中央选择性，崇祯期尤其需要奏疏、地方、朝鲜及满文材料交叉。 因而这里标示的是材料能够支持的行动与制度位置，而不是对动机、地方执行或普通人经验的无保留复原。

\eventanchor{ML-1596-001}{明·1596（万历二十四年）}
\paragraph{1596（万历二十四年）：万历皇帝派遣宦官担任税吏和矿监} 这一节点应放回本章已经展开的权力、财政、地方社会与边地关系中理解。账本把它出现的条件概括为：财政征解、交通工程或货币支付的压力使相关措施进入政务。 随后可追踪的变化是：其法定安排不等于地方执行，后续须以地域材料追踪负担与成效。 可由《神宗实录》万历二十四年条；《明会典》相关条目；地方志、黄册/契约/河工材料 互相核对；不过，译写候选以编年材料定位；实录记载反映朝廷选择，崇祯期须另以奏疏、地方及敌对政权材料交叉。 因而这里标示的是材料能够支持的行动与制度位置，而不是对动机、地方执行或普通人经验的无保留复原。

\eventanchor{ML-1596-002}{明·1596（万历二十四年）}
\paragraph{1596（万历二十四年）：朝鲜战争、海上贸易和朝贡使节的本年多方记录提供跨政权比较，外交语言不能替代地方经验。} 这一节点应放回本章已经展开的权力、财政、地方社会与边地关系中理解。账本把它出现的条件概括为：朝贡名义、边境安全与港口贸易的利益使交涉持续发生。 随后可追踪的变化是：它为后续册封、互市或海防安排留下文书与跨政权比较线索。 可由《神宗实录》万历二十四年条；《明史·外国传》；《朝鲜王朝实录》、琉球《历代宝案》或海外档案 互相核对；不过，桥接条目只标示可回查的年度问题与材料链；实录有中央选择性，崇祯期尤其需要奏疏、地方、朝鲜及满文材料交叉。 因而这里标示的是材料能够支持的行动与制度位置，而不是对动机、地方执行或普通人经验的无保留复原。

\eventanchor{ML-1596-003}{明·1596（万历二十四年）}
\paragraph{1596（万历二十四年）：书院、科举、出版与宗教活动的本年线索可连接知识网络，存世文本有阶层与地域偏差。} 这一节点应放回本章已经展开的权力、财政、地方社会与边地关系中理解。账本把它出现的条件概括为：制度、教育或知识传播的需要推动了相关行动或文本形成。 随后可追踪的变化是：它提供制度和传播史的锚点，不能据此推断社会整体接受。 可由《神宗实录》万历二十四年条；《明史·选举志》或《艺文志》；文集、碑刻或书目材料 互相核对；不过，桥接条目只标示可回查的年度问题与材料链；实录有中央选择性，崇祯期尤其需要奏疏、地方、朝鲜及满文材料交叉。 因而这里标示的是材料能够支持的行动与制度位置，而不是对动机、地方执行或普通人经验的无保留复原。

\eventanchor{ML-1596-004}{明·1596（万历二十四年）}
\paragraph{1596（万历二十四年）：考成、人事、立储和留中等本年材料标记皇帝、内阁与言官的协调关系，不可只用党争标签解释。} 这一节点应放回本章已经展开的权力、财政、地方社会与边地关系中理解。账本把它出现的条件概括为：继承、官僚协调或皇权运作的压力使问题进入朝廷程序。 随后可追踪的变化是：它影响后续人事与政策议程；动机和派系标签不能由单一叙事裁定。 可由《神宗实录》万历二十四年条；《明史·本纪》及相关列传；诏令、奏疏或会典版本 互相核对；不过，桥接条目只标示可回查的年度问题与材料链；实录有中央选择性，崇祯期尤其需要奏疏、地方、朝鲜及满文材料交叉。 因而这里标示的是材料能够支持的行动与制度位置，而不是对动机、地方执行或普通人经验的无保留复原。

\eventanchor{ML-1596-005}{明·1596（万历二十四年）}
\paragraph{1596（万历二十四年）：本年灾荒、河工和地方赈济材料可用于研究风险分配，灾害叙事和统计数字应区分。} 这一节点应放回本章已经展开的权力、财政、地方社会与边地关系中理解。账本把它出现的条件概括为：自然风险与地方报告促成朝廷或地方的应对记录。 随后可追踪的变化是：赈济、迁徙和损失的实际范围仍须以受灾地区材料分别核验。 可由《神宗实录》万历二十四年条；《明史·五行志》；受灾府县地方志与奏疏 互相核对；不过，桥接条目只标示可回查的年度问题与材料链；实录有中央选择性，崇祯期尤其需要奏疏、地方、朝鲜及满文材料交叉。 因而这里标示的是材料能够支持的行动与制度位置，而不是对动机、地方执行或普通人经验的无保留复原。

\eventanchor{ML-1597-001}{明·1597（万历二十五年）十月17日}
\paragraph{1597（万历二十五年）十月17日：稷山之战：日军向汉城进军被明军阻止} 这一节点应放回本章已经展开的权力、财政、地方社会与边地关系中理解。账本把它出现的条件概括为：朝贡名义、边境安全与港口贸易的利益使交涉持续发生。 随后可追踪的变化是：它为后续册封、互市或海防安排留下文书与跨政权比较线索。 可由《神宗实录》万历二十五年条；《明史·外国传》；《朝鲜王朝实录》、琉球《历代宝案》或海外档案 互相核对；不过，译写候选以编年材料定位；实录记载反映朝廷选择，崇祯期须另以奏疏、地方及敌对政权材料交叉。 因而这里标示的是材料能够支持的行动与制度位置，而不是对动机、地方执行或普通人经验的无保留复原。

\eventanchor{ML-1597-002}{明·1597（万历二十五年）}
\paragraph{1597（万历二十五年）：书院、科举、出版与宗教活动的本年线索可连接知识网络，存世文本有阶层与地域偏差。} 这一节点应放回本章已经展开的权力、财政、地方社会与边地关系中理解。账本把它出现的条件概括为：制度、教育或知识传播的需要推动了相关行动或文本形成。 随后可追踪的变化是：它提供制度和传播史的锚点，不能据此推断社会整体接受。 可由《神宗实录》万历二十五年条；《明史·选举志》或《艺文志》；文集、碑刻或书目材料 互相核对；不过，桥接条目只标示可回查的年度问题与材料链；实录有中央选择性，崇祯期尤其需要奏疏、地方、朝鲜及满文材料交叉。 因而这里标示的是材料能够支持的行动与制度位置，而不是对动机、地方执行或普通人经验的无保留复原。

\eventanchor{ML-1597-003}{明·1597（万历二十五年）}
\paragraph{1597（万历二十五年）：本年灾荒、河工和地方赈济材料可用于研究风险分配，灾害叙事和统计数字应区分。} 这一节点应放回本章已经展开的权力、财政、地方社会与边地关系中理解。账本把它出现的条件概括为：自然风险与地方报告促成朝廷或地方的应对记录。 随后可追踪的变化是：赈济、迁徙和损失的实际范围仍须以受灾地区材料分别核验。 可由《神宗实录》万历二十五年条；《明史·五行志》；受灾府县地方志与奏疏 互相核对；不过，桥接条目只标示可回查的年度问题与材料链；实录有中央选择性，崇祯期尤其需要奏疏、地方、朝鲜及满文材料交叉。 因而这里标示的是材料能够支持的行动与制度位置，而不是对动机、地方执行或普通人经验的无保留复原。

\eventanchor{ML-1597-004}{明·1597（万历二十五年）}
\paragraph{1597（万历二十五年）：辽东、西北、朝鲜或西南军务的本年调兵与战报记录跨区域动员，补给与兵额需要独立核验。} 这一节点应放回本章已经展开的权力、财政、地方社会与边地关系中理解。账本把它出现的条件概括为：前期边防、地方武装或跨区域资源调动的压力推动了该行动。 随后可追踪的变化是：它改变了后续调兵、补给或地方秩序的条件；战果和伤亡须分开核对。 可由《神宗实录》万历二十五年条；《明史·兵志》及相关列传；边镇、地方志或对方材料 互相核对；不过，桥接条目只标示可回查的年度问题与材料链；实录有中央选择性，崇祯期尤其需要奏疏、地方、朝鲜及满文材料交叉。 因而这里标示的是材料能够支持的行动与制度位置，而不是对动机、地方执行或普通人经验的无保留复原。

\eventanchor{ML-1597-005}{明·1597（万历二十五年）}
\paragraph{1597（万历二十五年）：清丈、一条鞭、漕运或军饷的本年文书构成万历财政的可核查节点，定额、报数和实收必须分列。} 这一节点应放回本章已经展开的权力、财政、地方社会与边地关系中理解。账本把它出现的条件概括为：财政征解、交通工程或货币支付的压力使相关措施进入政务。 随后可追踪的变化是：其法定安排不等于地方执行，后续须以地域材料追踪负担与成效。 可由《神宗实录》万历二十五年条；《明会典》相关条目；地方志、黄册/契约/河工材料 互相核对；不过，桥接条目只标示可回查的年度问题与材料链；实录有中央选择性，崇祯期尤其需要奏疏、地方、朝鲜及满文材料交叉。 因而这里标示的是材料能够支持的行动与制度位置，而不是对动机、地方执行或普通人经验的无保留复原。

\eventanchor{ML-1598-001}{明·1598（万历二十六年）一月4日}
\paragraph{1598（万历二十六年）一月4日：蔚山之围：明朝和朝鲜军队未能将日本人赶出蔚山城堡} 这一节点应放回本章已经展开的权力、财政、地方社会与边地关系中理解。账本把它出现的条件概括为：朝贡名义、边境安全与港口贸易的利益使交涉持续发生。 随后可追踪的变化是：它为后续册封、互市或海防安排留下文书与跨政权比较线索。 可由《神宗实录》万历二十六年条；《明史·外国传》；《朝鲜王朝实录》、琉球《历代宝案》或海外档案 互相核对；不过，译写候选以编年材料定位；实录记载反映朝廷选择，崇祯期须另以奏疏、地方及敌对政权材料交叉。 因而这里标示的是材料能够支持的行动与制度位置，而不是对动机、地方执行或普通人经验的无保留复原。

\eventanchor{ML-1598-002}{明·1598（万历二十六年）十月}
\paragraph{1598（万历二十六年）十月：泗川之战（1598）：明朝和朝鲜军队未能将日本人赶出泗川} 这一节点应放回本章已经展开的权力、财政、地方社会与边地关系中理解。账本把它出现的条件概括为：朝贡名义、边境安全与港口贸易的利益使交涉持续发生。 随后可追踪的变化是：它为后续册封、互市或海防安排留下文书与跨政权比较线索。 可由《神宗实录》万历二十六年条；《明史·外国传》；《朝鲜王朝实录》、琉球《历代宝案》或海外档案 互相核对；不过，译写候选以编年材料定位；实录记载反映朝廷选择，崇祯期须另以奏疏、地方及敌对政权材料交叉。 因而这里标示的是材料能够支持的行动与制度位置，而不是对动机、地方执行或普通人经验的无保留复原。

\eventanchor{ML-1598-003}{明·1598（万历二十六年）十月7日}
\paragraph{1598（万历二十六年）十月7日：围攻顺天：明朝和朝鲜军队未能将日本人逐出顺天城} 这一节点应放回本章已经展开的权力、财政、地方社会与边地关系中理解。账本把它出现的条件概括为：朝贡名义、边境安全与港口贸易的利益使交涉持续发生。 随后可追踪的变化是：它为后续册封、互市或海防安排留下文书与跨政权比较线索。 可由《神宗实录》万历二十六年条；《明史·外国传》；《朝鲜王朝实录》、琉球《历代宝案》或海外档案 互相核对；不过，译写候选以编年材料定位；实录记载反映朝廷选择，崇祯期须另以奏疏、地方及敌对政权材料交叉。 因而这里标示的是材料能够支持的行动与制度位置，而不是对动机、地方执行或普通人经验的无保留复原。

\subsection{1598年前后的材料线索}

\eventanchor{ML-1598-004}{明·1598（万历二十六年）十二月16日}
\paragraph{1598（万历二十六年）十二月16日：鹭梁海战：明朝鲜海军击败日本舰队} 这一节点应放回本章已经展开的权力、财政、地方社会与边地关系中理解。账本把它出现的条件概括为：朝贡名义、边境安全与港口贸易的利益使交涉持续发生。 随后可追踪的变化是：它为后续册封、互市或海防安排留下文书与跨政权比较线索。 可由《神宗实录》万历二十六年条；《明史·外国传》；《朝鲜王朝实录》、琉球《历代宝案》或海外档案 互相核对；不过，译写候选以编年材料定位；实录记载反映朝廷选择，崇祯期须另以奏疏、地方及敌对政权材料交叉。 因而这里标示的是材料能够支持的行动与制度位置，而不是对动机、地方执行或普通人经验的无保留复原。

\eventanchor{ML-1598-005}{明·1598（万历二十六年）十二月24日}
\paragraph{1598（万历二十六年）十二月24日：日本入侵朝鲜（1592-98）：日本军队从朝鲜撤军} 这一节点应放回本章已经展开的权力、财政、地方社会与边地关系中理解。账本把它出现的条件概括为：朝贡名义、边境安全与港口贸易的利益使交涉持续发生。 随后可追踪的变化是：它为后续册封、互市或海防安排留下文书与跨政权比较线索。 可由《神宗实录》万历二十六年条；《明史·外国传》；《朝鲜王朝实录》、琉球《历代宝案》或海外档案 互相核对；不过，译写候选以编年材料定位；实录记载反映朝廷选择，崇祯期须另以奏疏、地方及敌对政权材料交叉。 因而这里标示的是材料能够支持的行动与制度位置，而不是对动机、地方执行或普通人经验的无保留复原。

\eventanchor{ML-1598-006}{明·1598（万历二十六年）}
\paragraph{1598（万历二十六年）：亳州叛乱：苗族叛乱被镇压} 这一节点应放回本章已经展开的权力、财政、地方社会与边地关系中理解。账本把它出现的条件概括为：前期边防、地方武装或跨区域资源调动的压力推动了该行动。 随后可追踪的变化是：它改变了后续调兵、补给或地方秩序的条件；战果和伤亡须分开核对。 可由《神宗实录》万历二十六年条；《明史·兵志》及相关列传；边镇、地方志或对方材料 互相核对；不过，译写候选以编年材料定位；实录记载反映朝廷选择，崇祯期须另以奏疏、地方及敌对政权材料交叉。 因而这里标示的是材料能够支持的行动与制度位置，而不是对动机、地方执行或普通人经验的无保留复原。

\eventanchor{ML-1598-007}{明·1598（万历二十六年）}
\paragraph{1598（万历二十六年）：蒙古人杀明主帅李如松} 这一节点应放回本章已经展开的权力、财政、地方社会与边地关系中理解。账本把它出现的条件概括为：继承、官僚协调或皇权运作的压力使问题进入朝廷程序。 随后可追踪的变化是：它影响后续人事与政策议程；动机和派系标签不能由单一叙事裁定。 可由《神宗实录》万历二十六年条；《明史·本纪》及相关列传；诏令、奏疏或会典版本 互相核对；不过，译写候选以编年材料定位；实录记载反映朝廷选择，崇祯期须另以奏疏、地方及敌对政权材料交叉。 因而这里标示的是材料能够支持的行动与制度位置，而不是对动机、地方执行或普通人经验的无保留复原。

\eventanchor{ML-1598-008}{明·1598（万历二十六年）}
\paragraph{1598（万历二十六年）：明朝骑兵尝试发射三管火绳枪，然后用它作为盾牌，同时用另一只手用马刀攻击。} 这一节点应放回本章已经展开的权力、财政、地方社会与边地关系中理解。账本把它出现的条件概括为：继承、官僚协调或皇权运作的压力使问题进入朝廷程序。 随后可追踪的变化是：它影响后续人事与政策议程；动机和派系标签不能由单一叙事裁定。 可由《神宗实录》万历二十六年条；《明史·本纪》及相关列传；诏令、奏疏或会典版本 互相核对；不过，译写候选以编年材料定位；实录记载反映朝廷选择，崇祯期须另以奏疏、地方及敌对政权材料交叉。 因而这里标示的是材料能够支持的行动与制度位置，而不是对动机、地方执行或普通人经验的无保留复原。

\eventanchor{ML-1598-009}{明·1598（万历二十六年）}
\paragraph{1598（万历二十六年）：广东官员允许西班牙人在皮纳尔进行贸易} 这一节点应放回本章已经展开的权力、财政、地方社会与边地关系中理解。账本把它出现的条件概括为：朝贡名义、边境安全与港口贸易的利益使交涉持续发生。 随后可追踪的变化是：它为后续册封、互市或海防安排留下文书与跨政权比较线索。 可由《神宗实录》万历二十六年条；《明史·外国传》；《朝鲜王朝实录》、琉球《历代宝案》或海外档案 互相核对；不过，译写候选以编年材料定位；实录记载反映朝廷选择，崇祯期须另以奏疏、地方及敌对政权材料交叉。 因而这里标示的是材料能够支持的行动与制度位置，而不是对动机、地方执行或普通人经验的无保留复原。

\eventanchor{ML-1599-001}{明·1599（万历二十七年）}
\paragraph{1599（万历二十七年）：各大港口均有高级太监驻守} 这一节点应放回本章已经展开的权力、财政、地方社会与边地关系中理解。账本把它出现的条件概括为：继承、官僚协调或皇权运作的压力使问题进入朝廷程序。 随后可追踪的变化是：它影响后续人事与政策议程；动机和派系标签不能由单一叙事裁定。 可由《神宗实录》万历二十七年条；《明史·本纪》及相关列传；诏令、奏疏或会典版本 互相核对；不过，译写候选以编年材料定位；实录记载反映朝廷选择，崇祯期须另以奏疏、地方及敌对政权材料交叉。 因而这里标示的是材料能够支持的行动与制度位置，而不是对动机、地方执行或普通人经验的无保留复原。

\eventanchor{ML-1599-002}{明·1599（万历二十七年）}
\paragraph{1599（万历二十七年）：本年灾荒、河工和地方赈济材料可用于研究风险分配，灾害叙事和统计数字应区分。} 这一节点应放回本章已经展开的权力、财政、地方社会与边地关系中理解。账本把它出现的条件概括为：自然风险与地方报告促成朝廷或地方的应对记录。 随后可追踪的变化是：赈济、迁徙和损失的实际范围仍须以受灾地区材料分别核验。 可由《神宗实录》万历二十七年条；《明史·五行志》；受灾府县地方志与奏疏 互相核对；不过，桥接条目只标示可回查的年度问题与材料链；实录有中央选择性，崇祯期尤其需要奏疏、地方、朝鲜及满文材料交叉。 因而这里标示的是材料能够支持的行动与制度位置，而不是对动机、地方执行或普通人经验的无保留复原。

\eventanchor{ML-1599-003}{明·1599（万历二十七年）}
\paragraph{1599（万历二十七年）：清丈、一条鞭、漕运或军饷的本年文书构成万历财政的可核查节点，定额、报数和实收必须分列。} 这一节点应放回本章已经展开的权力、财政、地方社会与边地关系中理解。账本把它出现的条件概括为：财政征解、交通工程或货币支付的压力使相关措施进入政务。 随后可追踪的变化是：其法定安排不等于地方执行，后续须以地域材料追踪负担与成效。 可由《神宗实录》万历二十七年条；《明会典》相关条目；地方志、黄册/契约/河工材料 互相核对；不过，桥接条目只标示可回查的年度问题与材料链；实录有中央选择性，崇祯期尤其需要奏疏、地方、朝鲜及满文材料交叉。 因而这里标示的是材料能够支持的行动与制度位置，而不是对动机、地方执行或普通人经验的无保留复原。

\eventanchor{ML-1599-004}{明·1599（万历二十七年）}
\paragraph{1599（万历二十七年）：朝鲜战争、海上贸易和朝贡使节的本年多方记录提供跨政权比较，外交语言不能替代地方经验。} 这一节点应放回本章已经展开的权力、财政、地方社会与边地关系中理解。账本把它出现的条件概括为：朝贡名义、边境安全与港口贸易的利益使交涉持续发生。 随后可追踪的变化是：它为后续册封、互市或海防安排留下文书与跨政权比较线索。 可由《神宗实录》万历二十七年条；《明史·外国传》；《朝鲜王朝实录》、琉球《历代宝案》或海外档案 互相核对；不过，桥接条目只标示可回查的年度问题与材料链；实录有中央选择性，崇祯期尤其需要奏疏、地方、朝鲜及满文材料交叉。 因而这里标示的是材料能够支持的行动与制度位置，而不是对动机、地方执行或普通人经验的无保留复原。

\eventanchor{ML-1599-005}{明·1599（万历二十七年）}
\paragraph{1599（万历二十七年）：考成、人事、立储和留中等本年材料标记皇帝、内阁与言官的协调关系，不可只用党争标签解释。} 这一节点应放回本章已经展开的权力、财政、地方社会与边地关系中理解。账本把它出现的条件概括为：继承、官僚协调或皇权运作的压力使问题进入朝廷程序。 随后可追踪的变化是：它影响后续人事与政策议程；动机和派系标签不能由单一叙事裁定。 可由《神宗实录》万历二十七年条；《明史·本纪》及相关列传；诏令、奏疏或会典版本 互相核对；不过，桥接条目只标示可回查的年度问题与材料链；实录有中央选择性，崇祯期尤其需要奏疏、地方、朝鲜及满文材料交叉。 因而这里标示的是材料能够支持的行动与制度位置，而不是对动机、地方执行或普通人经验的无保留复原。

\eventanchor{ML-1600-001}{明·1600（万历二十八年）一月17}
\paragraph{1600（万历二十八年）一月17：澳门的葡萄牙人攻击兰帕考的西班牙人。西班牙人放弃埃尔皮纳尔。} 这一节点应放回本章已经展开的权力、财政、地方社会与边地关系中理解。账本把它出现的条件概括为：朝贡名义、边境安全与港口贸易的利益使交涉持续发生。 随后可追踪的变化是：它为后续册封、互市或海防安排留下文书与跨政权比较线索。 可由《神宗实录》万历二十八年条；《明史·外国传》；《朝鲜王朝实录》、琉球《历代宝案》或海外档案 互相核对；不过，译写候选以编年材料定位；实录记载反映朝廷选择，崇祯期须另以奏疏、地方及敌对政权材料交叉。 因而这里标示的是材料能够支持的行动与制度位置，而不是对动机、地方执行或普通人经验的无保留复原。

\eventanchor{ML-1600-002}{明·1600（万历二十八年）}
\paragraph{1600（万历二十八年）：欧洲藏书规模超过中国} 这一节点应放回本章已经展开的权力、财政、地方社会与边地关系中理解。账本把它出现的条件概括为：制度、教育或知识传播的需要推动了相关行动或文本形成。 随后可追踪的变化是：它提供制度和传播史的锚点，不能据此推断社会整体接受。 可由《神宗实录》万历二十八年条；《明史·选举志》或《艺文志》；文集、碑刻或书目材料 互相核对；不过，译写候选以编年材料定位；实录记载反映朝廷选择，崇祯期须另以奏疏、地方及敌对政权材料交叉。 因而这里标示的是材料能够支持的行动与制度位置，而不是对动机、地方执行或普通人经验的无保留复原。

\eventanchor{ML-1600-003}{明·1600（万历二十八年）}
\paragraph{1600（万历二十八年）：东林、书院、刻书和宗教活动的本年材料记录精英网络，社团存在不等于一致政治立场。} 这一节点应放回本章已经展开的权力、财政、地方社会与边地关系中理解。账本把它出现的条件概括为：制度、教育或知识传播的需要推动了相关行动或文本形成。 随后可追踪的变化是：它提供制度和传播史的锚点，不能据此推断社会整体接受。 可由《神宗实录》万历二十八年条；《明史·选举志》或《艺文志》；文集、碑刻或书目材料 互相核对；不过，桥接条目只标示可回查的年度问题与材料链；实录有中央选择性，崇祯期尤其需要奏疏、地方、朝鲜及满文材料交叉。 因而这里标示的是材料能够支持的行动与制度位置，而不是对动机、地方执行或普通人经验的无保留复原。

\eventanchor{ML-1600-004}{明·1600（万历二十八年）}
\paragraph{1600（万历二十八年）：本年灾害、疫病或地方赈济线索应以受灾地区文书和地方志复核，中央奏报不能覆盖全部社会。} 这一节点应放回本章已经展开的权力、财政、地方社会与边地关系中理解。账本把它出现的条件概括为：自然风险与地方报告促成朝廷或地方的应对记录。 随后可追踪的变化是：赈济、迁徙和损失的实际范围仍须以受灾地区材料分别核验。 可由《神宗实录》万历二十八年条；《明史·五行志》；受灾府县地方志与奏疏 互相核对；不过，桥接条目只标示可回查的年度问题与材料链；实录有中央选择性，崇祯期尤其需要奏疏、地方、朝鲜及满文材料交叉。 因而这里标示的是材料能够支持的行动与制度位置，而不是对动机、地方执行或普通人经验的无保留复原。

\eventanchor{ML-1600-005}{明·1600（万历二十八年）}
\paragraph{1600（万历二十八年）：后金兴起、朝鲜交涉和海上网络的本年多方材料揭示区域关系变化，政权名称与译名须按原文说明。} 这一节点应放回本章已经展开的权力、财政、地方社会与边地关系中理解。账本把它出现的条件概括为：朝贡名义、边境安全与港口贸易的利益使交涉持续发生。 随后可追踪的变化是：它为后续册封、互市或海防安排留下文书与跨政权比较线索。 可由《神宗实录》万历二十八年条；《明史·外国传》；《朝鲜王朝实录》、琉球《历代宝案》或海外档案 互相核对；不过，桥接条目只标示可回查的年度问题与材料链；实录有中央选择性，崇祯期尤其需要奏疏、地方、朝鲜及满文材料交叉。 因而这里标示的是材料能够支持的行动与制度位置，而不是对动机、地方执行或普通人经验的无保留复原。

\eventanchor{ML-1600-006}{明·1600（万历二十八年）}
\paragraph{1600（万历二十八年）：辽东战事、边镇防务和西南兵事的本年军令记录资源调动，战报数字应与朝鲜、满文或地方材料比照。} 这一节点应放回本章已经展开的权力、财政、地方社会与边地关系中理解。账本把它出现的条件概括为：前期边防、地方武装或跨区域资源调动的压力推动了该行动。 随后可追踪的变化是：它改变了后续调兵、补给或地方秩序的条件；战果和伤亡须分开核对。 可由《神宗实录》万历二十八年条；《明史·兵志》及相关列传；边镇、地方志或对方材料 互相核对；不过，桥接条目只标示可回查的年度问题与材料链；实录有中央选择性，崇祯期尤其需要奏疏、地方、朝鲜及满文材料交叉。 因而这里标示的是材料能够支持的行动与制度位置，而不是对动机、地方执行或普通人经验的无保留复原。

\eventanchor{ML-1601-001}{明·1601（万历二十九年）}
\paragraph{1601（万历二十九年）：万历末至天启初的立储、内阁和言路材料标记中枢运作的堵塞与调整，派系归类需回到具体文书。（材料核查重点：诏令或奏议的形成与发受链）} 这一节点应放回本章已经展开的权力、财政、地方社会与边地关系中理解。账本把它出现的条件概括为：继承、官僚协调或皇权运作的压力使问题进入朝廷程序。 随后可追踪的变化是：它影响后续人事与政策议程；动机和派系标签不能由单一叙事裁定。 可由《神宗实录》万历二十九年条；《明史·本纪》及相关列传；诏令、奏疏或会典版本 互相核对；不过，桥接条目只标示可回查的年度问题与材料链；实录有中央选择性，崇祯期尤其需要奏疏、地方、朝鲜及满文材料交叉。；本条特别核对诏令或奏议的形成与发受链。 因而这里标示的是材料能够支持的行动与制度位置，而不是对动机、地方执行或普通人经验的无保留复原。

\eventanchor{ML-1601-002}{明·1601（万历二十九年）}
\paragraph{1601（万历二十九年）：矿税、加派、漕运和辽饷的本年奏疏与账目提供财政压力线索，预算、欠额和实解不得混同。} 这一节点应放回本章已经展开的权力、财政、地方社会与边地关系中理解。账本把它出现的条件概括为：财政征解、交通工程或货币支付的压力使相关措施进入政务。 随后可追踪的变化是：其法定安排不等于地方执行，后续须以地域材料追踪负担与成效。 可由《神宗实录》万历二十九年条；《明会典》相关条目；地方志、黄册/契约/河工材料 互相核对；不过，桥接条目只标示可回查的年度问题与材料链；实录有中央选择性，崇祯期尤其需要奏疏、地方、朝鲜及满文材料交叉。 因而这里标示的是材料能够支持的行动与制度位置，而不是对动机、地方执行或普通人经验的无保留复原。

\eventanchor{ML-1601-003}{明·1601（万历二十九年）}
\paragraph{1601（万历二十九年）：本年灾害、疫病或地方赈济线索应以受灾地区文书和地方志复核，中央奏报不能覆盖全部社会。} 这一节点应放回本章已经展开的权力、财政、地方社会与边地关系中理解。账本把它出现的条件概括为：自然风险与地方报告促成朝廷或地方的应对记录。 随后可追踪的变化是：赈济、迁徙和损失的实际范围仍须以受灾地区材料分别核验。 可由《神宗实录》万历二十九年条；《明史·五行志》；受灾府县地方志与奏疏 互相核对；不过，桥接条目只标示可回查的年度问题与材料链；实录有中央选择性，崇祯期尤其需要奏疏、地方、朝鲜及满文材料交叉。 因而这里标示的是材料能够支持的行动与制度位置，而不是对动机、地方执行或普通人经验的无保留复原。

\subsection{1601年前后的材料线索}

\eventanchor{ML-1601-004}{明·1601（万历二十九年）}
\paragraph{1601（万历二十九年）：万历末至天启初的立储、内阁和言路材料标记中枢运作的堵塞与调整，派系归类需回到具体文书。（材料核查重点：报称信息的来源、抄传与行政分类）} 这一节点应放回本章已经展开的权力、财政、地方社会与边地关系中理解。账本把它出现的条件概括为：继承、官僚协调或皇权运作的压力使问题进入朝廷程序。 随后可追踪的变化是：它影响后续人事与政策议程；动机和派系标签不能由单一叙事裁定。 可由《神宗实录》万历二十九年条；《明史·本纪》及相关列传；诏令、奏疏或会典版本 互相核对；不过，桥接条目只标示可回查的年度问题与材料链；实录有中央选择性，崇祯期尤其需要奏疏、地方、朝鲜及满文材料交叉。；本条特别核对报称信息的来源、抄传与行政分类。 因而这里标示的是材料能够支持的行动与制度位置，而不是对动机、地方执行或普通人经验的无保留复原。

\eventanchor{ML-1601-005}{明·1601（万历二十九年）}
\paragraph{1601（万历二十九年）：东林、书院、刻书和宗教活动的本年材料记录精英网络，社团存在不等于一致政治立场。} 这一节点应放回本章已经展开的权力、财政、地方社会与边地关系中理解。账本把它出现的条件概括为：制度、教育或知识传播的需要推动了相关行动或文本形成。 随后可追踪的变化是：它提供制度和传播史的锚点，不能据此推断社会整体接受。 可由《神宗实录》万历二十九年条；《明史·选举志》或《艺文志》；文集、碑刻或书目材料 互相核对；不过，桥接条目只标示可回查的年度问题与材料链；实录有中央选择性，崇祯期尤其需要奏疏、地方、朝鲜及满文材料交叉。 因而这里标示的是材料能够支持的行动与制度位置，而不是对动机、地方执行或普通人经验的无保留复原。

\eventanchor{ML-1601-006}{明·1601（万历二十九年）}
\paragraph{1601（万历二十九年）：后金兴起、朝鲜交涉和海上网络的本年多方材料揭示区域关系变化，政权名称与译名须按原文说明。} 这一节点应放回本章已经展开的权力、财政、地方社会与边地关系中理解。账本把它出现的条件概括为：朝贡名义、边境安全与港口贸易的利益使交涉持续发生。 随后可追踪的变化是：它为后续册封、互市或海防安排留下文书与跨政权比较线索。 可由《神宗实录》万历二十九年条；《明史·外国传》；《朝鲜王朝实录》、琉球《历代宝案》或海外档案 互相核对；不过，桥接条目只标示可回查的年度问题与材料链；实录有中央选择性，崇祯期尤其需要奏疏、地方、朝鲜及满文材料交叉。 因而这里标示的是材料能够支持的行动与制度位置，而不是对动机、地方执行或普通人经验的无保留复原。

\eventanchor{ML-1602-001}{明·1602（万历三十年）}
\paragraph{1602（万历三十年）：利玛窦定居北京传播基督教} 这一节点应放回本章已经展开的权力、财政、地方社会与边地关系中理解。账本把它出现的条件概括为：继承、官僚协调或皇权运作的压力使问题进入朝廷程序。 随后可追踪的变化是：它影响后续人事与政策议程；动机和派系标签不能由单一叙事裁定。 可由《神宗实录》万历三十年条；《明史·本纪》及相关列传；诏令、奏疏或会典版本 互相核对；不过，译写候选以编年材料定位；实录记载反映朝廷选择，崇祯期须另以奏疏、地方及敌对政权材料交叉。 因而这里标示的是材料能够支持的行动与制度位置，而不是对动机、地方执行或普通人经验的无保留复原。

\eventanchor{ML-1602-002}{明·1602（万历三十年）}
\paragraph{1602（万历三十年）：辽东战事、边镇防务和西南兵事的本年军令记录资源调动，战报数字应与朝鲜、满文或地方材料比照。} 这一节点应放回本章已经展开的权力、财政、地方社会与边地关系中理解。账本把它出现的条件概括为：前期边防、地方武装或跨区域资源调动的压力推动了该行动。 随后可追踪的变化是：它改变了后续调兵、补给或地方秩序的条件；战果和伤亡须分开核对。 可由《神宗实录》万历三十年条；《明史·兵志》及相关列传；边镇、地方志或对方材料 互相核对；不过，桥接条目只标示可回查的年度问题与材料链；实录有中央选择性，崇祯期尤其需要奏疏、地方、朝鲜及满文材料交叉。 因而这里标示的是材料能够支持的行动与制度位置，而不是对动机、地方执行或普通人经验的无保留复原。

\eventanchor{ML-1602-003}{明·1602（万历三十年）}
\paragraph{1602（万历三十年）：万历末至天启初的立储、内阁和言路材料标记中枢运作的堵塞与调整，派系归类需回到具体文书。} 这一节点应放回本章已经展开的权力、财政、地方社会与边地关系中理解。账本把它出现的条件概括为：继承、官僚协调或皇权运作的压力使问题进入朝廷程序。 随后可追踪的变化是：它影响后续人事与政策议程；动机和派系标签不能由单一叙事裁定。 可由《神宗实录》万历三十年条；《明史·本纪》及相关列传；诏令、奏疏或会典版本 互相核对；不过，桥接条目只标示可回查的年度问题与材料链；实录有中央选择性，崇祯期尤其需要奏疏、地方、朝鲜及满文材料交叉。 因而这里标示的是材料能够支持的行动与制度位置，而不是对动机、地方执行或普通人经验的无保留复原。

\eventanchor{ML-1602-004}{明·1602（万历三十年）}
\paragraph{1602（万历三十年）：矿税、加派、漕运和辽饷的本年奏疏与账目提供财政压力线索，预算、欠额和实解不得混同。} 这一节点应放回本章已经展开的权力、财政、地方社会与边地关系中理解。账本把它出现的条件概括为：财政征解、交通工程或货币支付的压力使相关措施进入政务。 随后可追踪的变化是：其法定安排不等于地方执行，后续须以地域材料追踪负担与成效。 可由《神宗实录》万历三十年条；《明会典》相关条目；地方志、黄册/契约/河工材料 互相核对；不过，桥接条目只标示可回查的年度问题与材料链；实录有中央选择性，崇祯期尤其需要奏疏、地方、朝鲜及满文材料交叉。 因而这里标示的是材料能够支持的行动与制度位置，而不是对动机、地方执行或普通人经验的无保留复原。

\eventanchor{ML-1602-005}{明·1602（万历三十年）}
\paragraph{1602（万历三十年）：本年灾害、疫病或地方赈济线索应以受灾地区文书和地方志复核，中央奏报不能覆盖全部社会。} 这一节点应放回本章已经展开的权力、财政、地方社会与边地关系中理解。账本把它出现的条件概括为：自然风险与地方报告促成朝廷或地方的应对记录。 随后可追踪的变化是：赈济、迁徙和损失的实际范围仍须以受灾地区材料分别核验。 可由《神宗实录》万历三十年条；《明史·五行志》；受灾府县地方志与奏疏 互相核对；不过，桥接条目只标示可回查的年度问题与材料链；实录有中央选择性，崇祯期尤其需要奏疏、地方、朝鲜及满文材料交叉。 因而这里标示的是材料能够支持的行动与制度位置，而不是对动机、地方执行或普通人经验的无保留复原。

\eventanchor{ML-1602-006}{明·1602（万历三十年）}
\paragraph{1602（万历三十年）：东林、书院、刻书和宗教活动的本年材料记录精英网络，社团存在不等于一致政治立场。} 这一节点应放回本章已经展开的权力、财政、地方社会与边地关系中理解。账本把它出现的条件概括为：制度、教育或知识传播的需要推动了相关行动或文本形成。 随后可追踪的变化是：它提供制度和传播史的锚点，不能据此推断社会整体接受。 可由《神宗实录》万历三十年条；《明史·选举志》或《艺文志》；文集、碑刻或书目材料 互相核对；不过，桥接条目只标示可回查的年度问题与材料链；实录有中央选择性，崇祯期尤其需要奏疏、地方、朝鲜及满文材料交叉。 因而这里标示的是材料能够支持的行动与制度位置，而不是对动机、地方执行或普通人经验的无保留复原。

\eventanchor{ML-1603-001}{明·1603（万历三十一年）十月}
\paragraph{1603（万历三十一年）十月：桑利叛乱：西班牙人、日本人和菲律宾人屠杀马尼拉的华人；万历皇帝指责一名太监询问西班牙人是否可以在甲美地采矿，激怒了西班牙人} 这一节点应放回本章已经展开的权力、财政、地方社会与边地关系中理解。账本把它出现的条件概括为：朝贡名义、边境安全与港口贸易的利益使交涉持续发生。 随后可追踪的变化是：它为后续册封、互市或海防安排留下文书与跨政权比较线索。 可由《神宗实录》万历三十一年条；《明史·外国传》；《朝鲜王朝实录》、琉球《历代宝案》或海外档案 互相核对；不过，译写候选以编年材料定位；实录记载反映朝廷选择，崇祯期须另以奏疏、地方及敌对政权材料交叉。 因而这里标示的是材料能够支持的行动与制度位置，而不是对动机、地方执行或普通人经验的无保留复原。

\eventanchor{ML-1603-002}{明·1603（万历三十一年）}
\paragraph{1603（万历三十一年）：努尔哈赤与明朝将军同意划定领土边界} 这一节点应放回本章已经展开的权力、财政、地方社会与边地关系中理解。账本把它出现的条件概括为：继承、官僚协调或皇权运作的压力使问题进入朝廷程序。 随后可追踪的变化是：它影响后续人事与政策议程；动机和派系标签不能由单一叙事裁定。 可由《神宗实录》万历三十一年条；《明史·本纪》及相关列传；诏令、奏疏或会典版本 互相核对；不过，译写候选以编年材料定位；实录记载反映朝廷选择，崇祯期须另以奏疏、地方及敌对政权材料交叉。 因而这里标示的是材料能够支持的行动与制度位置，而不是对动机、地方执行或普通人经验的无保留复原。

\eventanchor{ML-1603-003}{明·1603（万历三十一年）}
\paragraph{1603（万历三十一年）：中国学者陈迪在明代反海盗任务期间在大佑安湾（台湾由此得名）呆了一段时间，首次对台湾原住民进行了重要描述} 这一节点应放回本章已经展开的权力、财政、地方社会与边地关系中理解。账本把它出现的条件概括为：前期边防、地方武装或跨区域资源调动的压力推动了该行动。 随后可追踪的变化是：它改变了后续调兵、补给或地方秩序的条件；战果和伤亡须分开核对。 可由《神宗实录》万历三十一年条；《明史·兵志》及相关列传；边镇、地方志或对方材料 互相核对；不过，译写候选以编年材料定位；实录记载反映朝廷选择，崇祯期须另以奏疏、地方及敌对政权材料交叉。 因而这里标示的是材料能够支持的行动与制度位置，而不是对动机、地方执行或普通人经验的无保留复原。

\eventanchor{ML-1603-004}{明·1603（万历三十一年）}
\paragraph{1603（万历三十一年）：辽东战事、边镇防务和西南兵事的本年军令记录资源调动，战报数字应与朝鲜、满文或地方材料比照。} 这一节点应放回本章已经展开的权力、财政、地方社会与边地关系中理解。账本把它出现的条件概括为：前期边防、地方武装或跨区域资源调动的压力推动了该行动。 随后可追踪的变化是：它改变了后续调兵、补给或地方秩序的条件；战果和伤亡须分开核对。 可由《神宗实录》万历三十一年条；《明史·兵志》及相关列传；边镇、地方志或对方材料 互相核对；不过，桥接条目只标示可回查的年度问题与材料链；实录有中央选择性，崇祯期尤其需要奏疏、地方、朝鲜及满文材料交叉。 因而这里标示的是材料能够支持的行动与制度位置，而不是对动机、地方执行或普通人经验的无保留复原。

\eventanchor{ML-1603-005}{明·1603（万历三十一年）}
\paragraph{1603（万历三十一年）：万历末至天启初的立储、内阁和言路材料标记中枢运作的堵塞与调整，派系归类需回到具体文书。} 这一节点应放回本章已经展开的权力、财政、地方社会与边地关系中理解。账本把它出现的条件概括为：继承、官僚协调或皇权运作的压力使问题进入朝廷程序。 随后可追踪的变化是：它影响后续人事与政策议程；动机和派系标签不能由单一叙事裁定。 可由《神宗实录》万历三十一年条；《明史·本纪》及相关列传；诏令、奏疏或会典版本 互相核对；不过，桥接条目只标示可回查的年度问题与材料链；实录有中央选择性，崇祯期尤其需要奏疏、地方、朝鲜及满文材料交叉。 因而这里标示的是材料能够支持的行动与制度位置，而不是对动机、地方执行或普通人经验的无保留复原。

\eventanchor{ML-1603-006}{明·1603（万历三十一年）}
\paragraph{1603（万历三十一年）：本年灾害、疫病或地方赈济线索应以受灾地区文书和地方志复核，中央奏报不能覆盖全部社会。} 这一节点应放回本章已经展开的权力、财政、地方社会与边地关系中理解。账本把它出现的条件概括为：自然风险与地方报告促成朝廷或地方的应对记录。 随后可追踪的变化是：赈济、迁徙和损失的实际范围仍须以受灾地区材料分别核验。 可由《神宗实录》万历三十一年条；《明史·五行志》；受灾府县地方志与奏疏 互相核对；不过，桥接条目只标示可回查的年度问题与材料链；实录有中央选择性，崇祯期尤其需要奏疏、地方、朝鲜及满文材料交叉。 因而这里标示的是材料能够支持的行动与制度位置，而不是对动机、地方执行或普通人经验的无保留复原。

\eventanchor{ML-1604-001}{明·1604（万历三十二年）}
\paragraph{1604（万历三十二年）：东林运动：东林书院成立} 这一节点应放回本章已经展开的权力、财政、地方社会与边地关系中理解。账本把它出现的条件概括为：制度、教育或知识传播的需要推动了相关行动或文本形成。 随后可追踪的变化是：它提供制度和传播史的锚点，不能据此推断社会整体接受。 可由《神宗实录》万历三十二年条；《明史·选举志》或《艺文志》；文集、碑刻或书目材料 互相核对；不过，译写候选以编年材料定位；实录记载反映朝廷选择，崇祯期须另以奏疏、地方及敌对政权材料交叉。 因而这里标示的是材料能够支持的行动与制度位置，而不是对动机、地方执行或普通人经验的无保留复原。

\eventanchor{ML-1604-002}{明·1604（万历三十二年）}
\paragraph{1604（万历三十二年）：后金兴起、朝鲜交涉和海上网络的本年多方材料揭示区域关系变化，政权名称与译名须按原文说明。（材料核查重点：诏令或奏议的形成与发受链）} 这一节点应放回本章已经展开的权力、财政、地方社会与边地关系中理解。账本把它出现的条件概括为：朝贡名义、边境安全与港口贸易的利益使交涉持续发生。 随后可追踪的变化是：它为后续册封、互市或海防安排留下文书与跨政权比较线索。 可由《神宗实录》万历三十二年条；《明史·外国传》；《朝鲜王朝实录》、琉球《历代宝案》或海外档案 互相核对；不过，桥接条目只标示可回查的年度问题与材料链；实录有中央选择性，崇祯期尤其需要奏疏、地方、朝鲜及满文材料交叉。；本条特别核对诏令或奏议的形成与发受链。 因而这里标示的是材料能够支持的行动与制度位置，而不是对动机、地方执行或普通人经验的无保留复原。

\eventanchor{ML-1604-003}{明·1604（万历三十二年）}
\paragraph{1604（万历三十二年）：矿税、加派、漕运和辽饷的本年奏疏与账目提供财政压力线索，预算、欠额和实解不得混同。（材料核查重点：诏令或奏议的形成与发受链）} 这一节点应放回本章已经展开的权力、财政、地方社会与边地关系中理解。账本把它出现的条件概括为：财政征解、交通工程或货币支付的压力使相关措施进入政务。 随后可追踪的变化是：其法定安排不等于地方执行，后续须以地域材料追踪负担与成效。 可由《神宗实录》万历三十二年条；《明会典》相关条目；地方志、黄册/契约/河工材料 互相核对；不过，桥接条目只标示可回查的年度问题与材料链；实录有中央选择性，崇祯期尤其需要奏疏、地方、朝鲜及满文材料交叉。；本条特别核对诏令或奏议的形成与发受链。 因而这里标示的是材料能够支持的行动与制度位置，而不是对动机、地方执行或普通人经验的无保留复原。

\eventanchor{ML-1604-004}{明·1604（万历三十二年）}
\paragraph{1604（万历三十二年）：后金兴起、朝鲜交涉和海上网络的本年多方材料揭示区域关系变化，政权名称与译名须按原文说明。（材料核查重点：报称信息的来源、抄传与行政分类）} 这一节点应放回本章已经展开的权力、财政、地方社会与边地关系中理解。账本把它出现的条件概括为：朝贡名义、边境安全与港口贸易的利益使交涉持续发生。 随后可追踪的变化是：它为后续册封、互市或海防安排留下文书与跨政权比较线索。 可由《神宗实录》万历三十二年条；《明史·外国传》；《朝鲜王朝实录》、琉球《历代宝案》或海外档案 互相核对；不过，桥接条目只标示可回查的年度问题与材料链；实录有中央选择性，崇祯期尤其需要奏疏、地方、朝鲜及满文材料交叉。；本条特别核对报称信息的来源、抄传与行政分类。 因而这里标示的是材料能够支持的行动与制度位置，而不是对动机、地方执行或普通人经验的无保留复原。

\eventanchor{ML-1604-005}{明·1604（万历三十二年）}
\paragraph{1604（万历三十二年）：矿税、加派、漕运和辽饷的本年奏疏与账目提供财政压力线索，预算、欠额和实解不得混同。（材料核查重点：报称信息的来源、抄传与行政分类）} 这一节点应放回本章已经展开的权力、财政、地方社会与边地关系中理解。账本把它出现的条件概括为：财政征解、交通工程或货币支付的压力使相关措施进入政务。 随后可追踪的变化是：其法定安排不等于地方执行，后续须以地域材料追踪负担与成效。 可由《神宗实录》万历三十二年条；《明会典》相关条目；地方志、黄册/契约/河工材料 互相核对；不过，桥接条目只标示可回查的年度问题与材料链；实录有中央选择性，崇祯期尤其需要奏疏、地方、朝鲜及满文材料交叉。；本条特别核对报称信息的来源、抄传与行政分类。 因而这里标示的是材料能够支持的行动与制度位置，而不是对动机、地方执行或普通人经验的无保留复原。

\subsection{1604年前后的材料线索}

\eventanchor{ML-1604-006}{明·1604（万历三十二年）}
\paragraph{1604（万历三十二年）：万历末至天启初的立储、内阁和言路材料标记中枢运作的堵塞与调整，派系归类需回到具体文书。} 这一节点应放回本章已经展开的权力、财政、地方社会与边地关系中理解。账本把它出现的条件概括为：继承、官僚协调或皇权运作的压力使问题进入朝廷程序。 随后可追踪的变化是：它影响后续人事与政策议程；动机和派系标签不能由单一叙事裁定。 可由《神宗实录》万历三十二年条；《明史·本纪》及相关列传；诏令、奏疏或会典版本 互相核对；不过，桥接条目只标示可回查的年度问题与材料链；实录有中央选择性，崇祯期尤其需要奏疏、地方、朝鲜及满文材料交叉。 因而这里标示的是材料能够支持的行动与制度位置，而不是对动机、地方执行或普通人经验的无保留复原。

\eventanchor{ML-1605-001}{明·1605（万历三十三年）六月}
\paragraph{1605（万历三十三年）六月：一道雷电击倒了天坛的旗杆，非常不吉利，导致一些官员辞职} 这一节点应放回本章已经展开的权力、财政、地方社会与边地关系中理解。账本把它出现的条件概括为：继承、官僚协调或皇权运作的压力使问题进入朝廷程序。 随后可追踪的变化是：它影响后续人事与政策议程；动机和派系标签不能由单一叙事裁定。 可由《神宗实录》万历三十三年条；《明史·本纪》及相关列传；诏令、奏疏或会典版本 互相核对；不过，译写候选以编年材料定位；实录记载反映朝廷选择，崇祯期须另以奏疏、地方及敌对政权材料交叉。 因而这里标示的是材料能够支持的行动与制度位置，而不是对动机、地方执行或普通人经验的无保留复原。

\eventanchor{ML-1605-002}{明·1605（万历三十三年）}
\paragraph{1605（万历三十三年）：东林、书院、刻书和宗教活动的本年材料记录精英网络，社团存在不等于一致政治立场。（材料核查重点：诏令或奏议的形成与发受链）} 这一节点应放回本章已经展开的权力、财政、地方社会与边地关系中理解。账本把它出现的条件概括为：制度、教育或知识传播的需要推动了相关行动或文本形成。 随后可追踪的变化是：它提供制度和传播史的锚点，不能据此推断社会整体接受。 可由《神宗实录》万历三十三年条；《明史·选举志》或《艺文志》；文集、碑刻或书目材料 互相核对；不过，桥接条目只标示可回查的年度问题与材料链；实录有中央选择性，崇祯期尤其需要奏疏、地方、朝鲜及满文材料交叉。；本条特别核对诏令或奏议的形成与发受链。 因而这里标示的是材料能够支持的行动与制度位置，而不是对动机、地方执行或普通人经验的无保留复原。

\eventanchor{ML-1605-003}{明·1605（万历三十三年）}
\paragraph{1605（万历三十三年）：辽东战事、边镇防务和西南兵事的本年军令记录资源调动，战报数字应与朝鲜、满文或地方材料比照。（材料核查重点：诏令或奏议的形成与发受链）} 这一节点应放回本章已经展开的权力、财政、地方社会与边地关系中理解。账本把它出现的条件概括为：前期边防、地方武装或跨区域资源调动的压力推动了该行动。 随后可追踪的变化是：它改变了后续调兵、补给或地方秩序的条件；战果和伤亡须分开核对。 可由《神宗实录》万历三十三年条；《明史·兵志》及相关列传；边镇、地方志或对方材料 互相核对；不过，桥接条目只标示可回查的年度问题与材料链；实录有中央选择性，崇祯期尤其需要奏疏、地方、朝鲜及满文材料交叉。；本条特别核对诏令或奏议的形成与发受链。 因而这里标示的是材料能够支持的行动与制度位置，而不是对动机、地方执行或普通人经验的无保留复原。

\eventanchor{ML-1605-004}{明·1605（万历三十三年）}
\paragraph{1605（万历三十三年）：东林、书院、刻书和宗教活动的本年材料记录精英网络，社团存在不等于一致政治立场。（材料核查重点：报称信息的来源、抄传与行政分类）} 这一节点应放回本章已经展开的权力、财政、地方社会与边地关系中理解。账本把它出现的条件概括为：制度、教育或知识传播的需要推动了相关行动或文本形成。 随后可追踪的变化是：它提供制度和传播史的锚点，不能据此推断社会整体接受。 可由《神宗实录》万历三十三年条；《明史·选举志》或《艺文志》；文集、碑刻或书目材料 互相核对；不过，桥接条目只标示可回查的年度问题与材料链；实录有中央选择性，崇祯期尤其需要奏疏、地方、朝鲜及满文材料交叉。；本条特别核对报称信息的来源、抄传与行政分类。 因而这里标示的是材料能够支持的行动与制度位置，而不是对动机、地方执行或普通人经验的无保留复原。

\eventanchor{ML-1605-005}{明·1605（万历三十三年）}
\paragraph{1605（万历三十三年）：辽东战事、边镇防务和西南兵事的本年军令记录资源调动，战报数字应与朝鲜、满文或地方材料比照。（材料核查重点：报称信息的来源、抄传与行政分类）} 这一节点应放回本章已经展开的权力、财政、地方社会与边地关系中理解。账本把它出现的条件概括为：前期边防、地方武装或跨区域资源调动的压力推动了该行动。 随后可追踪的变化是：它改变了后续调兵、补给或地方秩序的条件；战果和伤亡须分开核对。 可由《神宗实录》万历三十三年条；《明史·兵志》及相关列传；边镇、地方志或对方材料 互相核对；不过，桥接条目只标示可回查的年度问题与材料链；实录有中央选择性，崇祯期尤其需要奏疏、地方、朝鲜及满文材料交叉。；本条特别核对报称信息的来源、抄传与行政分类。 因而这里标示的是材料能够支持的行动与制度位置，而不是对动机、地方执行或普通人经验的无保留复原。

\eventanchor{ML-1605-006}{明·1605（万历三十三年）}
\paragraph{1605（万历三十三年）：矿税、加派、漕运和辽饷的本年奏疏与账目提供财政压力线索，预算、欠额和实解不得混同。} 这一节点应放回本章已经展开的权力、财政、地方社会与边地关系中理解。账本把它出现的条件概括为：财政征解、交通工程或货币支付的压力使相关措施进入政务。 随后可追踪的变化是：其法定安排不等于地方执行，后续须以地域材料追踪负担与成效。 可由《神宗实录》万历三十三年条；《明会典》相关条目；地方志、黄册/契约/河工材料 互相核对；不过，桥接条目只标示可回查的年度问题与材料链；实录有中央选择性，崇祯期尤其需要奏疏、地方、朝鲜及满文材料交叉。 因而这里标示的是材料能够支持的行动与制度位置，而不是对动机、地方执行或普通人经验的无保留复原。

\eventanchor{ML-1606-001}{明·1606（万历三十四年）}
\paragraph{1606（万历三十四年）：云南军官骚乱并杀害采矿太监杨荣} 这一节点应放回本章已经展开的权力、财政、地方社会与边地关系中理解。账本把它出现的条件概括为：继承、官僚协调或皇权运作的压力使问题进入朝廷程序。 随后可追踪的变化是：它影响后续人事与政策议程；动机和派系标签不能由单一叙事裁定。 可由《神宗实录》万历三十四年条；《明史·本纪》及相关列传；诏令、奏疏或会典版本 互相核对；不过，译写候选以编年材料定位；实录记载反映朝廷选择，崇祯期须另以奏疏、地方及敌对政权材料交叉。 因而这里标示的是材料能够支持的行动与制度位置，而不是对动机、地方执行或普通人经验的无保留复原。
